\subsection{Row-hook exactness}
\label{sec:hook-homology}

The Fields character has a finite principal-specialization form
that records both the exactness range and its first obstruction.
Write $(q;q)_m=\prod_{i=1}^m(1-q^i)$ and
$(q;q)_\infty=\prod_{i\ge1}(1-q^i)$.
Let $Z$ be the Frobenius alphabet for $S_n$ and set
$X=Z/(1-q)$ in plethystic notation. This means
$p_d[X]=p_d[Z]/(1-q^d)$.

Let $\mathscr H_n(q,u;Z)$ be the Frobenius series of the
$(1,1)$ bosonic--fermionic coinvariant ring: $q$ records polynomial
degree and $u$ exterior degree. Let $\mathscr E_n$ be the
moment Euler character after evaluation on a row superspace of
dimension $(1|1)$.

\begin{lemma}[Finite Fields identity and its first discrepancy]
\label{lem:fields-euler}
For $0\le b\le n$,
\begin{equation}\label{eq:finite-fields}
 [u^b]\mathscr H_n
  =\sum_{h=0}^b(-1)^{b-h}
       \frac{(q;q)_{n-h}}{(q;q)_{b-h}}e_h[X]h_{n-h}[X].
\end{equation}
The corresponding coefficient of $\mathscr E_n$ is obtained by
replacing $(q;q)_{n-h}$ with $(q;q)_\infty$. Hence the two
characters agree in every total degree at most $n$, and
\begin{equation}\label{eq:first-fields-discrepancy}
 [q^{n-b+1}u^b](\mathscr H_n-\mathscr E_n)
       =e_b[Z]h_{n-b}[Z]\qquad(0\le b\le n).
\end{equation}
\end{lemma}

\begin{proof}
The Fields theorem \cite[Theorem~5.2 and equation~(2.5)]{MRW} gives
\[
 \mathscr H_n
 =\sum_{\mu\vdash n}s_\mu[Z]
      \sum_{T\in\SYT(\mu)}
        q^{\maj(T)}
        \prod_{r=1}^{\des(T)}(1+uq^{-r}).
\]
Put
\[
 \mathcal A_n(t,q;Z)=
   \sum_{\mu\vdash n}s_\mu[Z]
       \sum_{T\in\SYT(\mu)}t^{\des(T)}q^{\maj(T)}.
\]
Principal specialization of Schur functions gives the
character-valued Eulerian identity
\[
 \frac{\mathcal A_n(t,q;Z)}
      {\prod_{i=0}^n(1-tq^i)}
  =\sum_{m\ge0}
       h_n\left[Z\frac{1-q^{m+1}}{1-q}\right]t^m.
\]
One obtains it by standardizing semistandard tableaux and summing
the weakly increasing entries with forced rises at descents.
The standardization convention using comajor index gives the same
identity, since tableau evacuation preserves the number of
descents and interchanges major and comajor indices.

Expand
\[
 h_n[X-q^{m+1}X]
   =\sum_{h=0}^n(-1)^h q^{(m+1)h}e_h[X]h_{n-h}[X].
\]
Multiplication by the denominator above and evaluation at
$t=q^{-h}$ leave just one summand, yielding
\[
 \mathcal A_n(q^{-h},q;Z)
 =q^{-\binom h2}(q;q)_h(q;q)_{n-h}e_h[X]h_{n-h}[X].
\]
The elementary-symmetric identity
\[
 e_b(q^{-1},\ldots,q^{-d})
  =\sum_{h=0}^b
       \frac{(-1)^{b-h}q^{\binom h2}}
            {(q;q)_h(q;q)_{b-h}}q^{-hd}
\]
follows from the finite $q$-binomial theorem. Apply it with
$d=\des(T)$ and use the preceding evaluation to obtain
\eqref{eq:finite-fields}.

The polynomial--exterior algebra has Frobenius series
\[
 \sum_{h=0}^n u^h e_h[X]h_{n-h}[X].
\]
The formal moment algebra contributes the reciprocal Euler factor
\[
 \frac{(q;q)_\infty}{\prod_{a\ge0}(1+uq^a)},\qquad
 \prod_{a\ge0}(1+uq^a)^{-1}
     =\sum_{j\ge0}\frac{(-u)^j}{(q;q)_j}.
\]
This proves the asserted formula for $\mathscr E_n$.

The difference $(q;q)_{n-h}-(q;q)_\infty$ begins in degree
$n-h+1$ with coefficient one. In fixed exterior degree $b$,
every term of the difference therefore has $q$-degree at
least $n-b+1$. At that first degree only $h=b$ contributes,
and the constant term of $e_b[X]h_{n-b}[X]$ is
$e_b[Z]h_{n-b}[Z]$. This proves both conclusions.
\end{proof}

\begin{theorem}[Row-hook moment acyclicity]
\label{thm:hook-acyclicity}
For every $n$, every row dimension $k$, every row hook
$\lambda\vdash D$, and every $D\le n$,
\[
 \Hom_{\GL_k}\!\left(
   \Schur_\lambda K^k,H_i(\mathcal K_{n,k,D})_D\right)=0
                  \qquad(i>0).
\]
Thus the full $S_n$ character in every row-hook sector is its
moment Euler character throughout $D\le n$.
\end{theorem}

\begin{proof}
Evaluate the polynomial row functors on $K^{1|1}$. The coefficient
algebra is
$P=K[x_1,\ldots,x_n]\otimes\Lambda(\theta_1,\ldots,\theta_n)$.
The free super moment algebra $S$ has even generators of
bidegree $(a,0)$ for $a\ge1$ and odd generators of bidegree
$(a,1)$ for $a\ge0$. They act by
$\sum_i x_i^a$ and $\sum_i\theta_i x_i^a$, respectively.
Its ordinary Hilbert series is
\[
 H_S(q,u)=\frac{\prod_{a\ge0}(1+uq^a)}{(q;q)_\infty}.
\]
Trace generation identifies $P/S_+P$ with the superspace
coinvariant ring. Lemma~\ref{lem:fields-euler}, after taking
dimensions, says exactly that
$H_P=H_SH_{P/S_+P}$ through total degree $n$.
Lemma~\ref{lem:truncated-freeness} forces positive Tor to vanish
there. The super Koszul resolution uses exterior powers on even
generators and symmetric powers on odd generators, so it is
precisely this evaluation of the moment complex.

To pass back to ordinary row spaces, express each homogeneous
polynomial functor as
$\bigoplus_\lambda\Schur_\lambda\otimes C_\lambda$.
In characteristic zero all maps act on the multiplicity spaces
$C_\lambda$, and this decomposition can be evaluated on a
super vector space with its signed tensor permutations.
It commutes with homology. A nonempty Schur functor is nonzero
on $K^{1|1}$ exactly when its partition is a hook, and
\[
 s_{(a,1^b)}[q\mid u]=(q+u)q^{a-1}u^b.
\]
Vanishing of the specialized homology therefore forces the
vanishing of every row-hook multiplicity, including its entire
$S_n$-module structure.
\end{proof}

